\section{Related Work on Hidden-State Probing for Routing}
\label{sec:related-work-routing}

\subsection{Hidden State Probing for Routing and Evaluation}
\label{subsec:hidden-state-probing}

Hidden-state probing has been applied to classification tasks. One line of this work has focused on exploiting hidden states to choose which general-purpose LLM, from a fixed pool, should answer a query. Results report that a probe trained on the hidden states can outperform a classifier trained on the query's surface text.

Varshney et al. trained a probe on prefill activations, the hidden states from the forward pass over the prompt before generation starts, to predict, for each candidate model, its probability of answering a query correctly. The probe outperformed baseline classifiers built from sentence embeddings, retrieval, and text-complexity features, reaching a mean \emph{AUROC} (area under the ROC curve, a threshold-independent measure of how well a classifier ranks positive cases above negative ones) of 0.856 against 0.804 for the strongest baseline on the frontier model pool \cite{varshneyLLMRouterRethinking2026}.

Lugoloobi et al. reported the same pattern for a model's own success on mathematics and coding tasks: a probe trained on activations read before generation outperformed a TF-IDF baseline, with the margin ranging from 0.05 to 0.20 AUROC across models and tasks \cite{lugoloobiLLMsEncodeTheir2026}.

A separate study by Wu, Wang et al. found that a router trained on cross-layer hidden states outperformed routers built on output probabilities, embeddings, and verbose self-evaluation, for choosing between a local and a cloud model, with relative improvements of 16.68\% in AUROC and 18.86\% on the high-accuracy deployment scenario over the best baseline \cite{wuFairComprehensiveEvaluation2026}.

Reading a decision from the hidden states before generation has also been applied to tasks beyond choosing which model should answer.

Chen et al. built two classifiers on hidden states read during prompt prefill, to classify a prompt's topic as math, code, or general text: a trained linear probe, and a training-free classifier that reduces the hidden state to summary statistics and compares them to a reference statistic computed for each candidate topic. Direct LLM inference on this same task was less reliable than the trained probe, particularly for smaller models under zero-shot prompting: its accuracy improved only with larger models and in-context learning, needed careful prompt design to reach a reasonable level, and failed outright for the smallest model tested. The training-free classifier trailed the LLMs' own decisions on several cases, but generalized better to out-of-distribution and ambiguous prompts \cite{chenFastIntentClassification2025}.

Cencerrado et al. trained a linear probe on activations read after a question but before generation, to predict whether a model's own forthcoming answer would be correct. On out-of-distribution questions, the probe generalized better than an embedding-based baseline and a baseline built from \emph{verbalized confidence} (a confidence score the model states in its own text output), though the embedding-based baseline matched or slightly exceeded the probe on in-distribution questions \cite{cencerradoNoAnswerNeeded2026}.

Healy et al. trained a classifier on an agent's final-layer hidden states, taken from the same forward pass that generated a tool call. The classifier flagged whether that call hallucinated the wrong function, malformed its parameters, or bypassed the available tools, reaching up to 86.4\% detection accuracy \cite{healyInternalRepresentationsIndicators2026}.

Sun et al. trained a linear probe on an LLM's hidden state at the last input token, read before generation, to predict whether a task needs a tool call at all. The probe reached an AUROC of 0.89 to 0.96 across six models, exceeding the same models' own verbalized reasoning about tool necessity \cite{sunLLMAgentsAlready2026}.

Across the works reviewed in this subsection, hidden states carry information about a decision beyond what the model's own output reveals. Because this information is available before generation, a probe reading it may supply a signal that can be compared against the model's eventual output. This motivates using a probe as an evaluation mechanism that flags a routing error before it takes effect.

\subsection{LatentGate}
\label{subsec:latentgate}

Ratnakar et al. came closest to routing among specialized destinations with a hidden-state probe. LatentGate is a router that reads hidden states from a frozen small language model and classifies a query directly against a pool of specialized agents, to reduce the latency and cost of a prompt-based LLM router \cite{ratnakarLatentGateLowLatencySemantic2026}. The paper reports that a prompt-based LLM router takes 1,500 to 2,000 milliseconds per query, with cost scaling with the number of agents, against about 28 milliseconds for LatentGate, and estimates a 99.98\% reduction in the monthly cost of routing 1 million queries \cite{ratnakarLatentGateLowLatencySemantic2026}. LatentGate reaches this lower cost because it stops the forward pass at a fixed interior layer of the small language model and never generates output tokens. A linear probe then classifies the pooled hidden state at that layer into one of the candidate agents. The frozen model never receives the list of candidate agents, and it never receives a routing instruction. The query alone is tokenized and passed through the model, under a maximum sequence length of 128 tokens \cite{ratnakarLatentGateLowLatencySemantic2026}, so that its hidden states serve as a representation of the query for the probe to classify.

LatentGate is evaluated on its own synthetic benchmark of 100 enterprise agents, built so that agents in the same domain share vocabulary on purpose. On this benchmark's out-of-distribution queries, paraphrased and deliberately built to overlap in vocabulary with other agents, LatentGate reaches 80.0\% accuracy, against 60.0\% for a sentence-embedding classifier, BGE-M3 \cite{ratnakarLatentGateLowLatencySemantic2026}. Because the model is never instructed to route, its own routing accuracy is never measured, so the probe is never compared against the model it reads from. LatentGate attributes about 65\% of its own errors on its benchmark to the deliberate cross-domain vocabulary overlap \cite{ratnakarLatentGateLowLatencySemantic2026}, but the paper never reports the accuracy of a bag-of-words classifier on this benchmark, so whether such a classifier fails for the same reason, and by how much, is not established. The only such baseline that LatentGate reports is on CLINC150, a separate 150-class intent-classification benchmark. There, LatentGate outperforms an SVM and TF-IDF classifier by 1.7 percentage points, 92.9\% against 91.2\% \cite{ratnakarLatentGateLowLatencySemantic2026}. The small accuracy difference suggests that most of the routing signal on that benchmark is already recoverable from surface text, so the probe's contribution is limited. LatentGate's scope is also limited to single-turn queries \cite{ratnakarLatentGateLowLatencySemantic2026}.

\subsection{What the Literature Leaves Open}
\label{subsec:literature-gap}

The hidden-state probing work reviewed in Section~\ref{subsec:hidden-state-probing} and Section~\ref{subsec:latentgate} evaluates every probe on single-turn queries. None of it tests whether a hidden-state probe still routes correctly once a query follows more than one turn of conversation. Furthermore, none of it compares a probe's routing accuracy against the model's own routing accuracy, the comparison that defines a probe's effectiveness as an evaluation mechanism.

As Section~\ref{subsec:latentgate} shows, LatentGate is the reviewed work closest in approach, but its purpose differs, because it produces the routing decision instead of checking a decision the coordinator has already made. It is also the only one that builds its benchmark to be hard for a sentence-embedding classifier, by giving agents in the same domain overlapping vocabulary on purpose, and then attributes its own errors to that overlap. It does not run the same check for a bag-of-words classifier. LatentGate's scale also differs from a coordinator's: it routes among 100 agents, not the roughly three to seven a production coordinator typically selects among, as Section~\ref{sec:multi-agent} shows. At that scale, a classifier can fail because there are many candidates to choose among, not because any two are hard to tell apart. LatentGate's benchmark does not establish whether the same difficulty holds when the candidate pool is small and the only source of difficulty is overlap between a few agents.
